\Needspace{9\baselineskip}\begin{center}\textit{Theorie der Wurzeln der Gleichung \(x^n-1=0\) (wobei \(n\) als Primzahl vorausgesetzt wird).} \textit{Lässt man die Wurzel \(1\) weg, so sind die übrigen \((\Omega)\) in der Gleichung \(X=x^{n-1}+x^{n-2}+\mathrm{usw.}+x+1=0\) enthalten.} \textit{Die Funktion \(X\) lässt sich nicht in Faktoren niedrigeren Grades zerlegen, deren Koeffizienten alle rational sind.}\end{center}


\Needspace{6\baselineskip}\Article{341}
\textsc{Satz.} \textit{Ist die Funktion \(X\) durch die Funktion niedrigeren Grades} \[P=x^\lambda+Ax^{\lambda-1}+Bx^{\lambda-2}+\ldots+Kx+L\] \textit{teilbar, so können die Koeffizienten \(A,\allowbreak B\ldots L\) nicht sämtlich ganze Zahlen sein.}

\textit{Beweis.} Es sei \(X=PQ\) und \(\mathfrak P\) die Gesamtheit der Wurzeln der Gleichung \(P=0\), \(\mathfrak Q\) die Gesamtheit der Wurzeln der Gleichung \(Q=0\), sodass \(\Omega\) aus \(\mathfrak P\) und \(\mathfrak Q\) zusammen besteht. Ferner sei \(\mathfrak R\) die Gesamtheit der zu den Wurzeln in \(\mathfrak P\) reziproken Wurzeln, \(\mathfrak S\) die Gesamtheit der zu den Wurzeln in \(\mathfrak Q\) reziproken Wurzeln. Die in \(\mathfrak R\) enthaltenen Wurzeln seien die Wurzeln der Gleichung \(R=0\) (die, wie man leicht sieht, \(x^\lambda+\frac KLx^{\lambda-1}+\mathrm{usw.}+\frac ALx+\frac1L=0\) lautet), die in \(\mathfrak S\) enthaltenen die Wurzeln der Gleichung \(S=0\). Offenbar bilden auch die Wurzeln \(\mathfrak R\) und \(\mathfrak S\) zusammen die Gesamtheit \(\Omega\), und es gilt \(RS=X\). Wir unterscheiden nun vier Fälle.
\SourceLeaf{423}

I. Wenn \(\mathfrak P\) mit \(\mathfrak R\) übereinstimmt und daher \(P=R\) gilt. In diesem Fall sind die Wurzeln in \(\mathfrak P\) offenbar stets paarweise reziprok, und somit ist \(P\) das Produkt aus \(\frac12\lambda\) quadratischen Faktoren der Form \(xx-2x\cos\omega+1\); da ein solcher Faktor \(=(x-\cos\omega)^2+\sin\omega^2\) ist, sieht man leicht, dass \(P\) für jeden reellen Wert von \(x\) notwendig einen positiven reellen Wert annimmt. Die Gleichungen, deren Wurzeln jeweils die Quadrate, Kuben, vierten Potenzen \(\ldots\) und Potenzen mit dem Exponenten \(n-1\) der Wurzeln in \(\mathfrak P\) sind, seien \(P^{\prime}=0\), \(P^{\prime\prime}=0,\allowbreak P^{\prime\prime\prime}=0,\allowbreak \ldots P^\nu=0\), und die Werte, welche die Funktionen \(P,\allowbreak P^{\prime},\allowbreak P^{\prime\prime}\ldots P^\nu\) bei Setzen von \(x=1\) annehmen, seien jeweils \(p,\allowbreak p^{\prime},\allowbreak p^{\prime\prime}\ldots p^\nu\). Dann ist nach dem Vorhergehenden \(p\) eine positive Größe, und aus ganz demselben Grund sind auch \(p^{\prime},\allowbreak p^{\prime\prime}\) usw. positiv. Nun ist \(p\) der Wert der Funktion \((1-t)(1-u)(1-v)\) usw., den sie annimmt, wenn man für \(t,\allowbreak u,\allowbreak v\) usw. die Wurzeln in \(\mathfrak P\) einsetzt; \(p^{\prime}\) ist der Wert derselben Funktion, wenn man für \(t,\allowbreak u,\allowbreak v\) usw. die Quadrate jener Wurzeln einsetzt, usw.; außerdem ist der Wert für \(t=1,\allowbreak u=1,\allowbreak v=1\) usw. offenbar \(=0\). Daher ist die Summe \(p+p^{\prime}+p^{\prime\prime}\ldots+p^\nu\) eine durch \(n\) teilbare ganze Zahl. Ferner sieht man leicht, dass das Produkt \(PP^{\prime} P^{\prime\prime}\ldots\) \(=X^\lambda\) ist und daher \(pp^{\prime} p^{\prime\prime}\ldots=n^\lambda\) gilt. Wären nun alle Koeffizienten in \(P\) rational, so wären nach Artikel 338 auch alle in \(P^{\prime},\allowbreak P^{\prime\prime}\) usw. rational; nach Artikel 42 müssten aber alle diese Koeffizienten ganze Zahlen sein. Daher wären auch \(p,\allowbreak p^{\prime},\allowbreak p^{\prime\prime}\) usw. sämtlich ganze Zahlen; da ihr Produkt \(n^\lambda\) ist und für ihre Anzahl \(n-1>\lambda\) gilt, müssten notwendig einige von ihnen (mindestens \(n-1-\lambda\)) gleich \(=1\), die übrigen aber gleich \(n\) oder einer Potenz von \(n\) sein. Sind also \(g\) von ihnen gleich \(=1\), so ist die Summe \(p+p^{\prime}+\mathrm{usw.}\) offenbar \(\equiv g\ (\mathrm{mod.}\ n)\) und somit sicher nicht durch \(n\) teilbar. Die Annahme kann daher nicht bestehen.

II. Wenn \(\mathfrak P\) und \(\mathfrak R\) zwar nicht übereinstimmen, aber einige gemeinsame Wurzeln enthalten, sei \(\mathfrak T\) deren Gesamtheit und \(T=0\) die Gleichung, deren Wurzeln sie sind. Dann ist \(T\) der größte gemeinsame Teiler der Funktionen \(P,\allowbreak R\) (wie aus der Theorie der Gleichungen bekannt ist). Die Wurzeln in \(\mathfrak T\) sind aber offenbar stets paarweise reziprok; daher können nach dem zuvor Bewiesenen die Koeffizienten in \(T\) nicht alle rational sein. Dies wären sie jedoch sicher, wenn alle in \(P\) und damit auch alle in \(R\) rational wären, wie unmittelbar aus der Art des Verfahrens zur Bestimmung des größten gemeinsamen Teilers folgt. Die Annahme ist daher absurd.

III. Wenn \(\mathfrak Q\) und \(\mathfrak S\) entweder übereinstimmen oder wenigstens gemeinsame Wurzeln enthalten, können auf genau dieselbe Weise die Koeffizienten in \(Q\) nicht alle rational sein; sie wären es aber, wenn alle in \(P\) rational wären. Dies ist also unmöglich.

IV. Haben dagegen weder \(\mathfrak P\) mit \(\mathfrak R\) noch \(\mathfrak Q\) mit \(\mathfrak S\) irgendeine gemeinsame\SourceLeaf{424}
Wurzel, so müssen alle Wurzeln \(\mathfrak P\) in \(\mathfrak S\) und alle \(\mathfrak Q\) in \(\mathfrak R\) enthalten sein, woraus \(P=S\) und \(Q=R\) folgt. Daher ist \(X=PQ\) das Produkt aus \(P\) und \(R\), das heißt \[\begin{gathered}
\text{aus}\quad x^\lambda+Ax^{\lambda-1}\ldots+Kx+L\\
\text{und}\quad x^\lambda+\frac KLx^{\lambda-1}\ldots+\frac ALx+\frac1L
\end{gathered}\] woraus durch Setzen von \(x=1\) \[nL=(1+A\ldots+K+L)^2\] folgt. Wären nun alle Koeffizienten in \(P\) rational und daher nach Artikel 42 auch ganze Zahlen, so müsste \(L\), das den letzten Koeffizienten in \(X\), nämlich die Einheit, teilen muss, notwendig \(=\pm1\) sein; daraus folgte, dass \(\pm n\) eine Quadratzahl wäre. Da dies der Voraussetzung widerspricht, kann die Annahme nicht bestehen.

Aus diesem Satz geht also hervor, dass, wie man auch \(X\) in Faktoren zerlegt, deren Koeffizienten wenigstens zum Teil irrational sind und sich daher nur durch eine Gleichung höheren Grades bestimmen lassen.

\Needspace{9\baselineskip}\begin{center}\textit{Das Ziel der folgenden Untersuchungen wird erläutert.}\end{center}


\Needspace{6\baselineskip}\Article{342}
Das Ziel der folgenden Untersuchungen, das kurz zu erläutern nützlich sein wird, besteht darin, \(X\) \textsc{schrittweise} in immer mehr Faktoren zu zerlegen und dabei ihre Koeffizienten durch Gleichungen möglichst niedrigen Grades zu bestimmen, bis man so zu Linearfaktoren oder zu den Wurzeln \(\Omega\) selbst gelangt. Wir werden nämlich zeigen: Zerlegt man die Zahl \(n-1\) auf beliebige Weise in ganzzahlige Faktoren \(\alpha,\allowbreak \beta,\allowbreak \gamma\) usw. (für die man jeweils Primzahlen nehmen kann), so lässt sich \(X\) in \(\alpha\) Faktoren vom Grad \(\frac{n-1}{\alpha}\) zerlegen, deren Koeffizienten durch eine Gleichung vom Grad \(\alpha\) bestimmt werden; jeder dieser Faktoren wiederum in \(\beta\) andere vom Grad \(\frac{n-1}{\alpha\beta}\) mithilfe einer Gleichung vom Grad \(\beta\) usw. Bezeichnet also \(\nu\) die Anzahl der Faktoren \(\alpha,\allowbreak \beta,\allowbreak \gamma\) usw., so führt sich die Bestimmung der Wurzeln \(\Omega\) auf die Lösung von \(\nu\) Gleichungen der Grade \(\alpha,\allowbreak \beta,\allowbreak \gamma\) usw. zurück. \textit{Zum Beispiel} sind für \(n=17\), wo \(n-1=2.2.2.2\) gilt, vier quadratische Gleichungen zu lösen; für \(n=73\) drei quadratische und zwei kubische.

Da im Folgenden sehr oft Potenzen der Wurzel \(r\) zu betrachten sind, deren Exponenten wiederum Potenzen sind, und solche Ausdrücke nur mit Mühe zu setzen sind, werden wir zur Erleichterung des Drucks künftig\SourceLeaf{425}
die folgende Abkürzung verwenden. Für \(r,\allowbreak rr,\allowbreak r^3\) usw. schreiben wir \([1],\allowbreak [2],\allowbreak [3]\) usw. und allgemein für \(r^\lambda\), wobei \(\lambda\) irgendeine ganze Zahl bezeichnet, \([\lambda]\). Solche Ausdrücke sind also noch nicht vollständig bestimmt, werden es aber, sobald für \(r\) oder \([1]\) eine bestimmte Wurzel aus \(\Omega\) gewählt wird. Allgemein sind somit \([\lambda],\allowbreak [\mu]\) gleich oder ungleich, je nachdem, ob \(\lambda,\allowbreak \mu\) nach dem Modul \(n\) kongruent oder inkongruent sind; ferner gilt \([0]=1;\allowbreak \ [\lambda].[\mu]=[\lambda+\mu];\allowbreak \ [\lambda]^\nu=[\lambda\nu]\); die Summe \([0]+[\lambda]+[2\lambda]\ldots+[(n-1)\lambda]\) ist entweder \(0\) oder \(n\), je nachdem, ob \(\lambda\) durch \(n\) nicht teilbar oder teilbar ist.

\Needspace{9\baselineskip}\begin{center}\textit{Sämtliche Wurzeln \(\Omega\) werden in gewisse Klassen (Perioden) eingeteilt.}\end{center}


\Needspace{6\baselineskip}\Article{343}
Ist für den Modul \(n\) die Zahl \(g\) eine Zahl derjenigen Art, die wir in Abschnitt III primitive Wurzel genannt haben, so sind die \(n-1\) Zahlen \(1,\allowbreak g,\allowbreak gg\ldots g^{n-2}\) den Zahlen \(1,\allowbreak 2,\allowbreak 3\ldots n-1\) nach dem Modul \(n\) kongruent, wenn auch in anderer Reihenfolge: Jede Zahl der einen Folge hat in der anderen eine kongruente Zahl. Daraus folgt unmittelbar, dass die Wurzeln \([1],\allowbreak [g],\allowbreak [gg]\ldots[g^{n-2}]\) mit \(\Omega\) übereinstimmen; ganz ebenso stimmen allgemeiner \[\relax[\lambda],[\lambda g],[\lambda gg]\ldots[\lambda g^{n-2}]\quad\text{mit}\quad\Omega\] überein, wobei \(\lambda\) irgendeine durch \(n\) nicht teilbare ganze Zahl bezeichnet. Da ferner \(g^{n-1}\equiv1\ (\mathrm{mod.}\ n)\) gilt, sieht man ohne Mühe, dass die beiden Wurzeln \([\lambda g^\mu],\allowbreak [\lambda g^\nu]\) identisch oder verschieden sind, je nachdem, ob \(\mu,\allowbreak \nu\) nach \(n-1\) kongruent oder inkongruent sind. Ist also \(G\) eine andere primitive Wurzel, so stimmen die Wurzeln \([1],\allowbreak [g]\ldots[g^{n-2}]\) auch mit \([1],\allowbreak [G]\ldots[G^{n-2}]\) überein, wenn man von der Reihenfolge absieht. Außerdem beweist man aber leicht: Ist \(e\) ein Teiler von \(n-1\) und setzt man \(n-1=ef,\allowbreak \ g^e=h\), \(G^e=H\), so sind auch die \(f\) Zahlen \(1,\allowbreak h,\allowbreak hh\ldots h^{f-1}\) den Zahlen \(1,\allowbreak H,\allowbreak H^2\ldots H^{f-1}\) nach \(n\) kongruent (ohne Rücksicht auf die Reihenfolge). Nehmen wir nämlich \(G\equiv g^\omega\ (\mathrm{mod.}\ n)\) an, und sei \(\mu\) eine beliebige positive Zahl \(<f\) und \(\nu\) der kleinste Rest von \(\mu\omega\ (\mathrm{mod.}\ f)\). Dann gilt \(\nu e\equiv\mu\omega e\ (\mathrm{mod.}\ n-1)\), daher \(g^{\nu e}\equiv g^{\mu\omega e}\equiv G^{\mu e}\ (\mathrm{mod.}\ n)\) oder \(H^\mu\equiv h^\nu\); das heißt, jede Zahl der letzteren Folge \(1,\allowbreak H,\allowbreak H^2\) usw. hat in der Folge \(1,\allowbreak h,\allowbreak hh\ldots\) eine kongruente Zahl, und ebenso umgekehrt. — Daraus ist klar, dass die \(f\) Wurzeln \([1],\allowbreak [h],\allowbreak [hh]\ldots[h^{f-1}]\) mit \([1],\allowbreak [H],\allowbreak [H^2]\ldots[H^{f-1}]\) identisch sind; allgemeiner sieht man auf dieselbe Weise leicht, dass \[\begin{gathered}
\relax[\lambda],[\lambda h],[\lambda hh]\ldots[\lambda h^{f-1}]\\
\text{mit}\quad[\lambda],[\lambda H],[\lambda H^2]\ldots[\lambda H^{f-1}]
\end{gathered}\] übereinstimmen. Die \textit{Summe} von solchen \(f\) Wurzeln \([\lambda]+[\lambda h]+\mathrm{usw.}+[\lambda h^{f-1}]\), die, da sie\SourceLeaf{426}
sich bei Wahl einer anderen primitiven Wurzel für \(g\) nicht ändert, als unabhängig von \(g\) anzusehen ist, bezeichnen wir mit \((f;\allowbreak \lambda)\); die \textit{Gesamtheit} derselben Wurzeln nennen wir die \textit{Periode} \((f,\allowbreak \lambda)\), wobei die Reihenfolge der Wurzeln unbeachtet bleibt\HistoricalNote{*)}{\normalfont Im Folgenden dürfen wir die Summe auch den Zahlenwert der Periode oder einfach die Periode nennen, wenn keine Mehrdeutigkeit zu befürchten ist.}. — Bei der Darstellung einer solchen Periode ist es zweckmäßig, die einzelnen Wurzeln, aus denen sie besteht, auf den einfachsten Ausdruck zurückzuführen, nämlich für die Zahlen \(\lambda,\allowbreak \lambda h,\allowbreak \lambda hh\) usw. die kleinsten Reste nach dem Modul \(n\) einzusetzen; nach deren Größe können, wenn man will, auch die Glieder der Periode geordnet werden.

\textit{Zum Beispiel} besteht für \(n=19\), wo \(2\) eine primitive Wurzel ist, die Periode \((6,\allowbreak 1)\) aus den Wurzeln \([1],\allowbreak [8],\allowbreak [64],\allowbreak [512],\allowbreak [4096],\allowbreak [32768]\) oder \([1],\allowbreak [7],\allowbreak [8],\allowbreak [11],\allowbreak [12],\allowbreak [18]\). Ebenso besteht die Periode \((6,\allowbreak 2)\) aus \([2],\allowbreak [3],\allowbreak [5],\allowbreak [14],\allowbreak [16],\allowbreak [17]\). Die Periode \((6,\allowbreak 3)\) ist mit der vorigen identisch. Die Periode \((6,\allowbreak 4)\) enthält \([4],\allowbreak [6],\allowbreak [9],\allowbreak [10],\allowbreak [13],\allowbreak [15]\).

\Needspace{9\baselineskip}\begin{center}\textit{Verschiedene Sätze über die Perioden der Wurzeln \(\Omega\).}\end{center}


\Needspace{6\baselineskip}\Article{344}
Bei diesen Perioden ergeben sich unmittelbar die folgenden Beobachtungen:

I. Da \(\lambda h^f\equiv\lambda,\allowbreak \ \lambda h^{f+1}\equiv\lambda h\) usw. \((\mathrm{mod.}\ n)\) gilt, bestehen offenbar aus denselben Wurzeln, aus denen \((f,\allowbreak \lambda)\) besteht, auch \((f,\allowbreak \lambda h),\allowbreak (f,\allowbreak \lambda hh)\) usw.; bezeichnet also allgemein \([\lambda']\) irgendeine Wurzel aus \((f,\allowbreak \lambda)\), so ist diese Periode mit \((f,\allowbreak \lambda')\) vollständig identisch. Haben also zwei Perioden mit gleich vielen Wurzeln (die wir \textit{gleichartig} nennen werden) irgendeine Wurzel gemeinsam, so sind sie offenbar identisch. Daher können nicht zwei Wurzeln gemeinsam in einer Periode enthalten sein, während in einer anderen gleichartigen nur eine von ihnen vorkommt. Ferner ist klar: Gehören zwei Wurzeln \([\lambda],\allowbreak [\lambda']\) zu derselben Periode mit \(f\) Gliedern, so ist der Wert des Ausdrucks \(\frac{\lambda'}{\lambda}\ (\mathrm{mod.}\ n)\) irgendeiner Potenz von \(h\) kongruent, oder man kann \(\lambda'\equiv\lambda g^{\nu e}\ (\mathrm{mod.}\ n)\) setzen.

II. Ist \(f=n-1,\allowbreak \ e=1\), so stimmt die Periode \((f,\allowbreak 1)\) offenbar mit \(\Omega\) überein; in den übrigen Fällen besteht \(\Omega\) aus \(e\) Perioden \((f,\allowbreak 1),\allowbreak (f,\allowbreak g),\allowbreak (f,\allowbreak gg)\ldots(f,\allowbreak g^{e-1})\). Diese Perioden sind also sämtlich voneinander verschieden, und offenbar stimmt jede andere gleichartige Periode \((f,\allowbreak \lambda)\) mit einer von ihnen überein, sofern \([\lambda]\) zu \(\Omega\) gehört, das heißt, wenn \(\lambda\) durch \(n\) nicht teilbar ist. Die Periode \((f,\allowbreak 0)\) oder \((f,\allowbreak kn)\) besteht dagegen offenbar aus \(f\) Einheiten. Ebenso leicht sieht man,\SourceLeaf{427}
dass, wenn \(\lambda\) irgendeine durch \(n\) nicht teilbare Zahl ist, auch die Gesamtheit der \(e\) Perioden \((f,\allowbreak \lambda),\allowbreak (f,\allowbreak \lambda g),\allowbreak (f,\allowbreak \lambda gg)\ldots(f,\allowbreak \lambda g^{e-1})\) mit \(\Omega\) übereinstimmt. — So besteht \textit{zum Beispiel} für \(n=19,\allowbreak \ f=6\) die Gesamtheit \(\Omega\) aus den drei Perioden \((6,\allowbreak 1),\allowbreak (6,\allowbreak 2),\allowbreak (6,\allowbreak 4)\), auf deren eine sich jede andere gleichartige außer \((6,\allowbreak 0)\) zurückführen lässt.

III. Ist \(n-1\) das Produkt aus drei positiven Zahlen \(a,\allowbreak b,\allowbreak c\), so besteht offenbar jede Periode mit \(bc\) Gliedern aus \(b\) Perioden mit \(c\) Gliedern, nämlich \((bc,\allowbreak \lambda)\) aus \((c,\allowbreak \lambda),\allowbreak (c,\allowbreak \lambda g^a),\allowbreak (c,\allowbreak \lambda g^{2a})\ldots(c,\allowbreak \lambda g^{ab-a})\); daher werden diese als in jener enthalten bezeichnet. So besteht für \(n=19\) die Periode \((6,\allowbreak 1)\) aus den drei Perioden \((2,\allowbreak 1),\allowbreak (2,\allowbreak 8),\allowbreak (2,\allowbreak 7)\), deren erste die Wurzeln \(r,\allowbreak r^{18}\), deren zweite \(r^8,\allowbreak r^{11}\) und deren dritte \(r^7,\allowbreak r^{12}\) enthält.

